\section{Otimização}\label{sec:otimizacao}

Neste trabalho, o dimensionamento das longarinas e das pranchas do tabuleiro é tratado como um problema de otimização de seção (\textit{sizing}) \parencite{BendsoeSigmund2003}, no qual a topologia do sistema estrutural, isto é, o número e a disposição das peças, é definida \textit{a priori} pelo projetista, e o algoritmo busca as dimensões ótimas dessas peças dentro dos limites estabelecidos.

\subsection{Histórico da otimização de pontes}\label{sec:historico_otimizacao_pontes}

A aplicação de técnicas de otimização ao projeto de pontes concentra-se, historicamente, em pontes de concreto e aço, antecedendo a própria difusão dos algoritmos evolutivos. Essa aplicação evoluiu de procedimentos determinísticos voltados à minimização de custo \parencite{Aguilar1973, Brar1984} para formulações que articulam múltiplos níveis de decisão, do componente à configuração estrutural, e múltiplos critérios de desempenho \parencite{CohnLounis1994, SimoesNegrao1994}.
A adoção de metaheurísticas ampliou essa formulação a problemas com numerosas variáveis discretas e verificações estruturais acopladas \parencite{Marti2013, GarciaSeguraYepes2016}, incorporando em seguida metas de confiabilidade ao longo da vida útil \parencite{GarciaSegura2017} e, mais recentemente, modelagem tridimensional acoplada a busca local \parencite{deMoraes2025}.

Para a madeira, os antecedentes concentram-se em passarelas estaiadas de madeira laminada colada, sistema estrutural distinto das pontes rodoviárias de madeira roliça estudadas neste trabalho. \textcite{Negrao1999} tratou a minimização de custo variando dimensões dos elementos e propriedades do material, e \textcite{SimoesNegrao2005} formulou um procedimento de projeto ótimo baseado em confiabilidade, integrando análise por elementos finitos, análise de sensibilidade e uma formulação multicritério de custo, tensões, deslocamentos e confiabilidade.
Esses antecedentes demonstram que tanto a busca de soluções econômicas quanto a consideração de incertezas já fazem parte da pesquisa sobre pontes de madeira, mas sem tratamento conjunto da variabilidade dimensional da madeira roliça e de uma formulação multiobjetivo, foco deste trabalho.

Esse tratamento de robustez, detalhado na Seção~\ref{sec:robustez}, avalia a resposta das soluções às variações prescritas nas dimensões; ele não corresponde, por si só, à imposição de uma probabilidade de falha ou de um índice-alvo de confiabilidade.
A formulação multiobjetivo utilizada é apresentada a seguir.

\subsection{Estrutura de otimização multiobjetivo}\label{sec:otim_multiobjetivo}

Um problema de otimização multiobjetivo envolve a minimização e/ou maximização simultânea de múltiplas funções objetivo, sujeitas a um conjunto de restrições que definem a viabilidade das soluções candidatas.
Diferentemente da otimização monoobjetivo, essa formulação não conduz a uma única solução ótima global.
Alternativamente, resulta em um conjunto de soluções ótimas de Pareto, cada uma representando um caso distinto de compromisso (\textit{trade-off}) entre critérios de desempenho concorrentes.
 
A formulação matemática geral de um problema de otimização multiobjetivo pode ser expressa conforme a Equação~\eqref{eq:mop}.
 
\begin{equation}
\left\{
\begin{array}{ll}
f_m(\mathbf{x}), & m = 1,2,\ldots,M; \\
g_j(\mathbf{x}) \leq 0, & j = 1,2,\ldots,J; \\
h_k(\mathbf{x}) = 0, & k = 1,2,\ldots,K; \\
x_i^{(L)} \leq x_i \leq x_i^{(U)}, & i = 1,2,\ldots,n.
\end{array}
\right.
\label{eq:mop}
\end{equation}
 
 Nesta formulação, o vetor de funções objetivo é definido segundo a Equação
 ~\eqref{eq:fx}.
 
 \begin{equation}
 f(\mathbf{x}) = \left(f_1(\mathbf{x}), f_2(\mathbf{x}), \ldots, f_M(\mathbf{x})\right)^T\label{eq:fx}
 \end{equation}
 
 Onde $\mathbf{x} = (x_1, x_2, \ldots, x_n)^T$ 
representa o vetor $n$-dimensional de variáveis de projeto.
O espaço de projeto viável é delimitado pelas restrições de desigualdade $g_j(\mathbf{x})$, restrições de igualdade $h_k(\mathbf{x})$ e restrições de domínio definidas pelos limites inferior $x_i^{(L)}$ e superior $x_i^{(U)}$, que em conjunto estabelecem o domínio admissível $(D)$ \parencite{DEB2009}.

Objetivos inicialmente concebidos como problemas de maximização são convertidos em problemas equivalentes de minimização mediante a multiplicação das respectivas funções objetivo por $-1$, garantindo uma estrutura computacional unificada.

 \subsection{Algoritmo de otimização NSGA-II} \label{sec:nsga2}
 
 Desenvolvido por \textcite{DEB2009}, este algoritmo evolutivo consolidou-se na literatura por sua capacidade de aproximar de forma eficiente o conjunto ótimo de Pareto. 
O NSGA-II preserva a natureza vetorial do problema multiobjetivo e busca um conjunto diverso de soluções não dominadas, sem recorrer à escalarização (soma ponderada) das funções objetivo.
O mecanismo central do algoritmo baseia-se na classificação rápida de não dominância (\textit{Fast Non-Dominated Sorting}).
A cada geração iterativa, toda a população de indivíduos é avaliada e ordenada em frentes sucessivas de não dominância $\mathcal{F}_1, \mathcal{F}_2, \ldots, \mathcal{F}_k$.
A primeira frente ($\mathcal{F}_1$) agrupa as soluções temporariamente ótimas daquela geração, constituindo a aproximação corrente da fronteira de Pareto.
A frente subsequente ($\mathcal{F}_2$) abriga soluções que seriam dominantes caso os indivíduos da frente $\mathcal{F}_1$ fossem desconsiderados, e assim por diante.Para garantir que a busca não convirja prematuramente para uma única região e mantenha uma dispersão uniforme ao longo da fronteira, o algoritmo utiliza o operador de distância de aglomeração (\textit{crowding distance}).
Esta métrica estima a densidade de soluções em torno de um determinado indivíduo, calculando o hipervolume formado pelos seus vizinhos mais próximos.
Para um problema com $M$ funções objetivo, a distância de aglomeração $d_i$ do indivíduo $i$ é calculada conforme a Equação~\eqref{eq:crowding}.
 
 \begin{equation}d_i = \sum_{m=1}^{M} \frac{f_m(\mathbf{x}_{i+1}) - f_m(\mathbf{x}_{i-1})}{f_m^{\max} - f_m^{\min}}
 \label{eq:crowding}
 \end{equation}
 
O NSGA-II emprega uma estratégia fortemente elitista.
A cada ciclo, as populações parental e descendente são combinadas e ordenadas primeiramente pelo nível de Pareto e, em caso de empate na mesma frente, pela maior distância de aglomeração.
O Algoritmo~\ref{alg:nsga2} resume a lógica evolutiva e o fluxo de avaliações estruturais, incorporando a proteção contra incertezas descrita na Seção~\ref{sec:robustez}.
 
 \begin{algorithm}[!ht]
 \caption{Otimização robusta multiobjetivo via NSGA-II}
 \label{alg:nsga2}
 \begin{algorithmic}[1]
 \Require Limites $[x_i^{(L)}, x_i^{(U)}]$, tamanho da população $N_{pop}$, gerações $N_{gen}$, desvio de robustez $\rho$, grade de robustez $\Lambda = \{-1,\,-\tfrac{1}{2},\,0,\,\tfrac{1}{2},\,1\}$
 \Ensure Fronteira eficiente $\mathcal{F}_1$
 \State Inicializar população $P_0$ por amostragem aleatória em $[x_i^{(L)}, x_i^{(U)}]$
 \For{cada geração $t = 0$ até $N_{gen}-1$}
   \For{cada indivíduo $\mathbf{x} = (d, b_w, h, n_{long}, n_{tab}) \in P_t$}
     \State $(\mathbf{F}, \cdot) \gets$ \Call{AvaliarModelo}{$\mathbf{x}$} \Comment{objetivos no ponto nominal, sem agregação}
     \State $\mathbf{G} \gets -\boldsymbol{\infty}$
     \For{cada $\lambda \in \Lambda$} \Comment{grade de robustez, aplicada só a $d$}
       \State $\tilde{d} \gets d\,(1 + \rho\,\lambda)$
       \State $\tilde{\mathbf{x}} \gets (\tilde{d},\,b_w,\,h,\,n_{long},\,n_{tab})$
       \State $(\cdot, \mathbf{g}) \gets$ \Call{AvaliarModelo}{$\tilde{\mathbf{x}}$}
       \State $\mathbf{G} \gets \max(\mathbf{G}, \mathbf{g})$ \Comment{componente a componente, pior caso}
     \EndFor
   \EndFor
   \State $Q_t \gets$ seleção por torneio, cruzamento SBX e mutação polinomial sobre $P_t$
   \State $R_t \gets P_t \cup Q_t$
   \State Ordenar $R_t$ em frentes não dominadas $\{\mathcal{F}_1, \mathcal{F}_2, \ldots\}$
   \State $P_{t+1} \gets$ melhores $N_{pop}$ indivíduos por nível de Pareto e distância de aglomeração
 \EndFor
 \State \Return $\mathcal{F}_1$
 \end{algorithmic}
 \end{algorithm}

Para a evolução das gerações, o algoritmo baseia-se em operadores genéticos específicos para o tratamento de variáveis contínuas.
A aplicação conjunta do cruzamento simulado binário (SBX) e da mutação polinomial (PM) permite refinar as soluções locais no espaço de busca, ao passo que a eliminação sistemática de duplicatas resguarda a diversidade genética da população, prevenindo a convergência prematura.
Os parâmetros de configuração adotados em todas as execuções deste trabalho estão reunidos na Tabela~\ref{tab:nsga2}.

\begin{table}[!ht]
\caption{Parâmetros de configuração do NSGA-II.\label{tab:nsga2}}
\begin{tabular*}{\columnwidth}{@{\extracolsep\fill}lc@{\extracolsep\fill}}
\toprule
Parâmetro & Valor \\
\midrule
Tamanho da população, $N_{pop}$ & 50 \\
Número de gerações, $N_{gen}$ & 300 \\
Amostragem inicial & Aleatória uniforme \\
Cruzamento & SBX, $p_c = 0{,}90$, $\eta_c = 15$ \\
Mutação & Polinomial, $\eta_m = 20$ \\
Eliminação de duplicatas & Sim \\
Semente aleatória & 1 \\
Número de pontos da grade de robustez, $N_c$ & 5 \\
Desvio de robustez, $\rho$ & 5\% \\
\bottomrule
\end{tabular*}
\end{table}
 
 \subsection{Robustez na otimização}\label{sec:robustez}
 
Na engenharia estrutural, a otimização puramente determinística tende a posicionar as soluções ótimas exatamente sobre a fronteira matemática das restrições ativas.
Essa característica as torna extremamente sensíveis a pequenas variações nas variáveis de projeto.
Na prática, essa sensibilidade é agravada pela variabilidade inerente às propriedades físicas dos materiais empregados e pelas tolerâncias naturais de fabricação, corte e montagem das peças.
Sob condições reais, uma estrutura cujo projeto encontre-se no limite exato de viabilidade determinística possui alta probabilidade de violar normas de segurança ou estados limites de serviço caso suas dimensões finais desviem minimamente do previsto.
Para mitigar esse efeito, o conceito de robustez deve ser incorporado de forma nativa ao processo algorítmico.
A estratégia adotada consiste em manter as funções objetivo no valor nominal de projeto, mas avaliar a viabilidade de cada solução candidata também sob a influência de um conjunto de perturbações controladas.

A perturbação é aplicada exclusivamente ao diâmetro da longarina, $d$. Tal condição é explicada pelo fato de que a madeira roliça possui variabilidade natural de diâmetro relevante para o comportamento estrutural. Em vez de amostrar $\xi$ aleatoriamente, a perturbação segue uma grade determinística de cinco pontos, $\Lambda = \{-1,\,-\tfrac{1}{2},\,0,\,\tfrac{1}{2},\,1\}$: essa escolha elimina o viés de semente que uma amostragem aleatória de poucos pontos introduziria, ao mesmo tempo em que os cinco níveis bastam para capturar eventuais degraus de viabilidade dentro do intervalo $[-\rho,\,\rho]$, sem depender de a resposta ser monotônica.
Para um dado indivíduo $\mathbf{x} = (d, b_w, h, n_{long}, n_{tab})$, são geradas $N_c = |\Lambda| = 5$ realizações perturbadas, denotadas por $\tilde{\mathbf{x}}^{(c)}$, conforme definido na Equação~\eqref{eq:perturbacao}.

 \begin{equation}\tilde{d}^{(c)} = d\left(1 + \rho\,\lambda^{(c)}\right),\qquad \lambda^{(c)} \in \Lambda,\qquad \tilde{x}_i^{(c)} = x_i \ \text{para}\ i \neq d\label{eq:perturbacao}
 \end{equation}

O termo $\rho$ representa a magnitude do desvio paramétrico tolerado pelo projeto em relação ao valor nominal do diâmetro, e $\lambda^{(c)}$ percorre os cinco pontos da grade $\Lambda$.
Quando $\rho = 0$, a grade colapsa para o único ponto $\lambda = 0$ e o procedimento se reduz à avaliação determinística padrão ($N_c = 1$).
Em cada ciclo, as avaliações perturbadas são processadas independentemente no modelo analítico, porém apenas as restrições são reavaliadas nesses pontos: as funções objetivo permanecem calculadas no ponto nominal $\mathbf{x}$, sem qualquer agregação sobre a grade.
Posteriormente, os valores das restrições nas $N_c$ realizações perturbadas são agregados pelo operador de pior caso (máximo), de acordo com a Equação~\eqref{eq:gmed}.

 \begin{equation}g_j(\mathbf{x}) = \max_{c=1,\ldots,N_c}\ g_j\!\left(\tilde{\mathbf{x}}^{(c)}\right)
 \label{eq:gmed}
 \end{equation}

Por meio desta formulação, uma solução só é julgada admissível se todas as suas restrições permanecerem viáveis ($g_j \leq 0$) em cada ponto da grade de perturbação do diâmetro, ou seja, no pior caso: projetos configurados muito próximos das restrições violam ao menos um dos pontos desfavoráveis da grade, o que já basta para que sejam penalizados pelo mecanismo de viabilidade do algoritmo.
O resultado é o deslocamento natural da fronteira de Pareto rumo a zonas mais conservadoras do espaço viável de projeto.

